% Options for packages loaded elsewhere
\PassOptionsToPackage{unicode}{hyperref}
\PassOptionsToPackage{hyphens}{url}
\PassOptionsToPackage{dvipsnames,svgnames,x11names}{xcolor}
%
\documentclass[
  letterpaper,
  DIV=11,
  numbers=noendperiod]{scrartcl}

\usepackage{amsmath,amssymb}
\usepackage{iftex}
\ifPDFTeX
  \usepackage[T1]{fontenc}
  \usepackage[utf8]{inputenc}
  \usepackage{textcomp} % provide euro and other symbols
\else % if luatex or xetex
  \usepackage{unicode-math}
  \defaultfontfeatures{Scale=MatchLowercase}
  \defaultfontfeatures[\rmfamily]{Ligatures=TeX,Scale=1}
\fi
\usepackage{lmodern}
\ifPDFTeX\else  
    % xetex/luatex font selection
\fi
% Use upquote if available, for straight quotes in verbatim environments
\IfFileExists{upquote.sty}{\usepackage{upquote}}{}
\IfFileExists{microtype.sty}{% use microtype if available
  \usepackage[]{microtype}
  \UseMicrotypeSet[protrusion]{basicmath} % disable protrusion for tt fonts
}{}
\makeatletter
\@ifundefined{KOMAClassName}{% if non-KOMA class
  \IfFileExists{parskip.sty}{%
    \usepackage{parskip}
  }{% else
    \setlength{\parindent}{0pt}
    \setlength{\parskip}{6pt plus 2pt minus 1pt}}
}{% if KOMA class
  \KOMAoptions{parskip=half}}
\makeatother
\usepackage{xcolor}
\setlength{\emergencystretch}{3em} % prevent overfull lines
\setcounter{secnumdepth}{5}
% Make \paragraph and \subparagraph free-standing
\makeatletter
\ifx\paragraph\undefined\else
  \let\oldparagraph\paragraph
  \renewcommand{\paragraph}{
    \@ifstar
      \xxxParagraphStar
      \xxxParagraphNoStar
  }
  \newcommand{\xxxParagraphStar}[1]{\oldparagraph*{#1}\mbox{}}
  \newcommand{\xxxParagraphNoStar}[1]{\oldparagraph{#1}\mbox{}}
\fi
\ifx\subparagraph\undefined\else
  \let\oldsubparagraph\subparagraph
  \renewcommand{\subparagraph}{
    \@ifstar
      \xxxSubParagraphStar
      \xxxSubParagraphNoStar
  }
  \newcommand{\xxxSubParagraphStar}[1]{\oldsubparagraph*{#1}\mbox{}}
  \newcommand{\xxxSubParagraphNoStar}[1]{\oldsubparagraph{#1}\mbox{}}
\fi
\makeatother


\providecommand{\tightlist}{%
  \setlength{\itemsep}{0pt}\setlength{\parskip}{0pt}}\usepackage{longtable,booktabs,array}
\usepackage{calc} % for calculating minipage widths
% Correct order of tables after \paragraph or \subparagraph
\usepackage{etoolbox}
\makeatletter
\patchcmd\longtable{\par}{\if@noskipsec\mbox{}\fi\par}{}{}
\makeatother
% Allow footnotes in longtable head/foot
\IfFileExists{footnotehyper.sty}{\usepackage{footnotehyper}}{\usepackage{footnote}}
\makesavenoteenv{longtable}
\usepackage{graphicx}
\makeatletter
\newsavebox\pandoc@box
\newcommand*\pandocbounded[1]{% scales image to fit in text height/width
  \sbox\pandoc@box{#1}%
  \Gscale@div\@tempa{\textheight}{\dimexpr\ht\pandoc@box+\dp\pandoc@box\relax}%
  \Gscale@div\@tempb{\linewidth}{\wd\pandoc@box}%
  \ifdim\@tempb\p@<\@tempa\p@\let\@tempa\@tempb\fi% select the smaller of both
  \ifdim\@tempa\p@<\p@\scalebox{\@tempa}{\usebox\pandoc@box}%
  \else\usebox{\pandoc@box}%
  \fi%
}
% Set default figure placement to htbp
\def\fps@figure{htbp}
\makeatother

\KOMAoption{captions}{tableheading}
\makeatletter
\@ifpackageloaded{caption}{}{\usepackage{caption}}
\AtBeginDocument{%
\ifdefined\contentsname
  \renewcommand*\contentsname{Table of contents}
\else
  \newcommand\contentsname{Table of contents}
\fi
\ifdefined\listfigurename
  \renewcommand*\listfigurename{List of Figures}
\else
  \newcommand\listfigurename{List of Figures}
\fi
\ifdefined\listtablename
  \renewcommand*\listtablename{List of Tables}
\else
  \newcommand\listtablename{List of Tables}
\fi
\ifdefined\figurename
  \renewcommand*\figurename{Figure}
\else
  \newcommand\figurename{Figure}
\fi
\ifdefined\tablename
  \renewcommand*\tablename{Table}
\else
  \newcommand\tablename{Table}
\fi
}
\@ifpackageloaded{float}{}{\usepackage{float}}
\floatstyle{ruled}
\@ifundefined{c@chapter}{\newfloat{codelisting}{h}{lop}}{\newfloat{codelisting}{h}{lop}[chapter]}
\floatname{codelisting}{Listing}
\newcommand*\listoflistings{\listof{codelisting}{List of Listings}}
\makeatother
\makeatletter
\makeatother
\makeatletter
\@ifpackageloaded{caption}{}{\usepackage{caption}}
\@ifpackageloaded{subcaption}{}{\usepackage{subcaption}}
\makeatother

\usepackage{bookmark}

\IfFileExists{xurl.sty}{\usepackage{xurl}}{} % add URL line breaks if available
\urlstyle{same} % disable monospaced font for URLs
\hypersetup{
  pdftitle={300\_tech\_notes.md},
  colorlinks=true,
  linkcolor={blue},
  filecolor={Maroon},
  citecolor={Blue},
  urlcolor={Blue},
  pdfcreator={LaTeX via pandoc}}


\title{300\_tech\_notes.md}
\author{}
\date{}

\begin{document}
\maketitle

\renewcommand*\contentsname{Table of contents}
{
\hypersetup{linkcolor=}
\setcounter{tocdepth}{3}
\tableofcontents
}

LLM NOTES? go here: 320\_LLM\_agents\_NOTES.md

\section{300 Tech Notes}\label{tech-notes}

typst file

\subsection{Jul 2026}\label{jul-2026}

\begin{itemize}
\tightlist
\item
  Change from \emph{.typ file to }.md; why? it simple text and md easier
\end{itemize}

\subsubsection{sudo apt}\label{sudo-apt}

sudo apt autoremove \# a dependency no longer needed

sudo apt-mark \# pkgs settings, such as being held

sudo apt install --fix-broken sudo apt install --fix-missing

\subsection{linux cli}\label{linux-cli}

df -h \# free space on system

journalctl -b -1 \# show system messages from prior boot

\subsubsection{Audio}\label{audio}

WirePlumber PipeWire ALSA (bottom)

\begin{itemize}
\tightlist
\item
  Audio: low-level Pipewire, replaces pulseaudio
  (https://wiki.archlinux.org/title/PipeWire) (ALSA remains the lowest
  level library)
\item
  Audio: \emph{WirePlumper} (wpctl) is used to manage this library
\item
  Audio: May see references to older PulseAudio, but this probably
  PipeWire emulating the older package. PulseAudo package is no longer
  on machine.
\item
  \emph{wpctl} is CLI
\end{itemize}

** Mar 2026

\subsubsection{Use Calibre for all files - one central place - including
those currently on
iPad.}\label{use-calibre-for-all-files---one-central-place---including-those-currently-on-ipad.}

From iPad: can upload to Drive! To iPad: from Drive to iPad? - from
Linux? iPad runs ``Calibre Companion'' to retrieve from linux - NOT the
reverse - From iCloud to linux? - sftp between linux \& ipad? (linux:
sudo systemctl start ssh; then ftp localhost)

** Feb 2026 Haskell working hard on data science; tutorials\ldots{}
cloud based GUI

** Dec 2025 - HP boot reports amdgpu psp gfx failed, but continues -
Status: unresolved - HP 845 G8, laptop; otherwise runs fine. - sudo
dmesg - sudo journalctr -b - secureddisplay? LOAD\_TA(0x1)

\subsubsection{Hardware}\label{hardware}

23 DEC 2024 - wireless mouse, \$10, bestbuy, M220 logitech, 18 month
batter life, plain vanilla is just fine.

\subsubsection{folds}\label{folds}

folds: use za or (define-key evil-normal-state-map (kbd ``'')
'evil-toggle-fold) need package for evil/folds - za open fold - zc close
fold - zr open all folds - zm close all folds - C-3 S-TAB (show outline
upto level 3 - WORKS!)

\subsubsection{url links}\label{url-links}

nytimes.com @ nytimes file:280\_emacs\_notes.qmd @ emacs notes
https://www.nytimes.com @ nytimes.com

** TODO - use folds, to simplify navigating this document - no R here,
no Stats - includes \emph{i3}, \emph{emacs}, \emph{Debian} - HTTR2 notes
= mess - Break this document into separate files, merge into one big
one?

\begin{itemize}
\tightlist
\item
  \textless2024-10-10 Thu\textgreater{} now Debian 12 (Cinnamon or i3);
  no more linux mint; LMDE \ldots{} \textless2024-08-26
  Mon\textgreater{}
\item
  do not use emacs as ``DESKTOP'' - everything, windows on the screen at
  once.
\item
  humans can not multitask; keep window only used for current work.
\item
  computers great for keeping lists of buffers, search when you move to
  next task.
\item
  tools like ivy can recall things like split window, code 1 side; dir
  on other
\item
  exist tools to maintain buffers in `projects'
\item
  EXWM still not sure.
\end{itemize}

\textless2024-07-01 Mon\textgreater{} - ex looks very useful! - SO(3) -
not sure what I am asking; normal subgroups, symmetries?

** Github actions - do not guess, do not play: READ DOCS - render quarto
(in gfm, github format, SEE \emph{gh\_actions\_project/quarto, SEE
main.yaml} - will render qmd to md file (github format) - do not soft
link qmd files, actual file or copy must in project - of course, can
render to md locally and merely push to github; - do not know branch and
yaml the same (ie not both main\ldots)

\begin{itemize}
\tightlist
\item
  (Github actions) Reading

  \begin{itemize}
  \tightlist
  \item
    Book \url{https://orchid00.github.io/actions_sandbox/}
  \item
    \url{https://www.r-bloggers.com/2021/01/automatic-rendering-of-a-plot-with-github-actions-3/}
  \item
    \url{https://www.simonpcouch.com/blog/2020-12-27-r-github-actions-commit/index.html}
  \item
    \url{https://fromthebottomoftheheap.net/2020/04/30/rendering-your-readme-with-github-actions/}
  \end{itemize}
\end{itemize}

** Posit Connect Cloud (new) - works with Shiny; gets code and manifest
(json, with all R dependencies) from github (works)

** 29-May-2024: Upgrade to R 4.4, Linux Mint 21.3 (Virginia)

\textless2024-09-25 Wed\textgreater{}

** ``ex'' editor

\begin{itemize}
\tightlist
\item
  simple, vi-like editor (see
  https://www.geeksforgeeks.org/ex-command-in-linux-with-examples/)
\item
  :a (to add text)
\item
  :w , :wq
\item
  print : 1,3 or :1,3p :4 (print line 4)
\item
  delete : 1,2d
\end{itemize}

prior to `ex', was `ed'

** PGP - install public key: \emph{gpg --recv-keys XXX}, where XXX= 16
alphanumber - KEY SIGNING: (summary) I can create public-private
keypair. If I `sign' doc with private key, then anyone with my public
key can READ and will know ONLY I could have encryted the file. -
Encrypt with my PUBLIC KEY, then only someone with my PRIVATE KEY can
read it. - ALL keys belong on Debian key ring. *** R - This is a
nightmare. There are several discrete steps; Do one at time! - sketch:
get public key; import; convert .asc to .gpg (if needed); cp .gpg to
Debian's ring - USE gpg; DO NOT USE apt-key - apt-secure - ONLY 2 line
to run. Just do what Cran R says to do for Debian, key-ring:
https://ftp.osuosl.org/pub/cran/. Then reboot, update apt. - gpg keys;
Debian: see apt-secure, READ carefully, but it doeswork. - gpg
--list-keys

** i3, i3wm Videos: basic config:
https://www.youtube.com/watch?v=88o2XuH3E08\&list=PLbcglKxZP5PMTWVe1KhZX3Jvo0xqk8h5I\&index=1
- man i3 has good help, but quick becomes \emph{x11} - review man
i3-dmenu-desktop

** PURPOSE: Misc Tech Notes; details, notes can be here (but COMMANDS
put on INDEX C). Math, R, probability and stats notes \textbf{do NOT go
here}

** emacs/org mode/eLisp

CAUTION: org mode is structured outline; markdown, typst are for
document formatting. TODO: Which to use?

SEE: Separate emacs/org documentation\\
281\_emacs\_notes.typ

SEE: 000\_PROJECTS.org

** functional programming (julia, rust, lisp) - Julia easier? - RUST
learning curve, but high performance - Both have functional?, but seems
to learn functional stay with lisp, racket, haskell or any similar

** IMPERATIVE v DECLARATIVE - math notation is DECLARATIVE; indicates
meaning, not exact details (example1: summation sign). Tell computer
what you WANT, leave details to compiler or program. - IMPERATIVE
(older, obsolete?) every detail, machine language?

** QUARTO: 2023-12-06 - Try couple of cv or resume templates and one
quarto extension. Seemed to be more work than worth. Just use Rmarkdown
to create resume. 2024-01-26 READ r4ds Ch 28-29 b/c Quarto and config
knitr

To render a Quarto (*.qmd) file with embeded typst. Use regular quarto
render methods.

** \{X11, Wayland, video\} \textbf{Warning:} Stay with X11 (proven,
well-supported, works)

Wayland is open source \emph{protocol} replacement for X windows,
(widgets says get graphics from X or Wayland). WESTON is reference
implementation.

Because Wayland accepts same toolkits (Qt, GTK) as X, the impact for
developers and users should be minimal. Wayland shrinks X; much
functional now in Linux kernel.

\begin{itemize}
\tightlist
\item
  Xserver: `display' server, runs locally, makes display and keyboard
  available to apps (either local or network)
  \textasciitilde/.xsession-errors - Unlike cli, GUIs have no console;
  errors err redirected to this file \textasciitilde/.Xauthority - is
  random code to control which progams output to my X11 session
\end{itemize}

** \{PDF, Pandoc, Latex\}

\subsubsection{PDF Notes {[}ignores html, css; also ignores YAML header
(pandoc \&
::render(){]}}\label{pdf-notes-ignores-html-css-also-ignores-yaml-header-pandoc-render}

\subsubsection{Print raw text (example, from github,
raw)}\label{print-raw-text-example-from-github-raw}

\emph{Quarto} saveas qmd, add yaml header, enclose all text in ```
(verbatim); render inside rstudio.

2023-12-30 - tlmgr controls much latex install: fonts, *.sty, - for PDF:
pdf2latex, pdflatex (pandoc will do it, but complain) - in R,
tinytex::latexmk()

To create pdf, just about everything works: pandoc, markdown, latex,
knitr.. Note: zathura uses library, open source \texttt{popler}.

NOTE: Missing latex .sty ?\\
- With .tex file, run (in R) tinytex:latemk(*.tex) to install - OR, use
tlmgr install - fonts installed? fc-list : family - also luaotfload
(loads fonts) - ENGINE=software (such as \textbf{luatex}, tex, pdftex) -
FORMAT=macros (such as \textbf{lualatex})

(Jan 2022) \textbf{Missing font, package? TinyTex} * update R *
keep\_tex: true (in YAML) * at R console tinytex::lualatex(``\ldots.
.tex''), or tinytex::latexmk(``\emph{.tex'') } \textasciitilde{} some
times works, sometimes not \textasciitilde{}

Lua in \emph{.tex file } see \textasciitilde/code/publish\_project/TEX/

(pre-Quarto !) HTML {[}to produce HTML with pandoc, all latex is
IGNORED.{]}

I do \textbf{not} know how to create fancy HTML files from knitr,
pandoc.

HTML is pain in ass and HUGE time waste. Pandoc can handle markdown and
small amounts of latex (math) b/c ppl have added filters or other
widgets to pandoc.

If using Latex, its packages, diagrams with Latex \ldots{} must go with
PDF.

-H header\\
-V or --variable\\
--pdf-engine=xelatex

\begin{itemize}
\tightlist
\item
  Try verbatum; process as a markdown. pandoc balks at processing
  straight text if it thinks it sees markdown. If lucky, !pandoc \% -o
  file.pdf will work.
\end{itemize}

\subsubsection{section\{LATEX NOTES\}}\label{sectionlatex-notes}

\begin{itemize}
\item
  Tikz seems to be most popular way to gaphics vs \textbf{pstricks}.

  footnote:

  \par

  \textbf{postscript} a more powerful programming language than tex;
  \textbf{pstricks}, ghostscript; pdf (a subset of postscript) hails
  from this. However, using postscript with latex requires addins, such
  as ghostscript; drivers; \ldots{} Avoid \textbf{postscript} and
  packages pstricks, even if greater capability.
\end{itemize}

Original tex was 320 low-level cmds (aka primitives). Macros created
from these. But actual engine (tex) hidden from user.

\textbf{LuaTex} (engine) is re-write of core TEX engine (hard, written
in C).\\
Therefore, \textbf{LuaTex} added primitives, more open (can use tex or
lua)

\textbf{LuaLatex} is macro package.

EXAMPLE: In .tex file, write lua: directlua is new primitive; lua api
inside value for

\[\pi = \directlua{tex.sprint(math.pi)}\]

\subsubsection{revealjs slides with
Quarto}\label{revealjs-slides-with-quarto}

** GIT commands

HEAD - normally points to branch (tip) but can point to commit (detached
HEAD). \textbf{Goes where you go, like a shadow.} Commits are immutable.
HEADS can move around.

\begin{itemize}
\item
  git merge --no-ff --no-commit main (no fastforward, \emph{dry-run})
\item
  git restore --source master\textasciitilde2 path/file (2 commits back,
  give path to file)
\item
  git log MANY USEFUL! See Book 2.3
  (https://git-scm.com/book/en/v2/Git-Basics-Viewing-the-Commit-History),
  Refernce (https://git-scm.com/docs/git-log)
\item
  git diff HEAD\textasciitilde{} HEAD \#show change between prior and
  most recent commit (``https://git-scm.com/docs/git-diff\_) book:
  (https://git-scm.com/book/en/v2/Git-Basics-Recording-Changes-to-the-Repository.html\#\_git\_diff\_staged)
\item
  git diff HEAD -- path/to/file.R \# this file only, show diff (Ex: git
  diff play dev \# merge back into dev, but 1st see changes)
\item[$\square$]
  UNDO one or more commits (many ways) Detach head by git checkout ,
  then build out new branch (git checkout -b new) git revert \# keeps
  history, undo one past commit git revert HEAD\textasciitilde2..HEAD \#
  keep history, but undo last 2 commits
\item
  git reflog \# how HEAD jumped around
\item
  git checkout HEAD@\{n\} \# move back to n ** Install R

  \begin{itemize}
  \tightlist
  \item
    on Ubuntu, or mint linux virgina, use jammy and follow
    https://cran.r-project.org/bin/linux/ubuntu/ (works)\\
  \item
    on Debian (such as LMDE) follow
    https://cran.r-project.org/bin/linux/debian/

    \begin{enumerate}
    \def\labelenumi{(\arabic{enumi})}
    \tightlist
    \item
      Add: sudo vim.tiny /etc/apt/sources.list
    \item
      insert: the deb XXX link
    \item
      save, run as jim, sudo apt install r-base etc.
    \end{enumerate}
  \end{itemize}
\end{itemize}

** REGEX - TODO import (?) all REGEX/* files to here

\begin{itemize}
\tightlist
\item
  for regex reading see 300\_tech\_reading.md
\end{itemize}

\section{-----------------------}\label{section}

\subsection{DOCUMENT REGEX HERE}\label{document-regex-here}

\subsection{(text, no examples in this
file)}\label{text-no-examples-in-this-file}

\section{-----------------------}\label{section-1}

/home/jim/code/docs/tech\_notes/REGEX.md

2024-04-30 - Use a cheat sheet - Reduce paper - Annotate (here, or in 2
files) ONLY when needs

\subsection{this file:
\textasciitilde/code/docs/tech\_notes/REGEX.md}\label{this-file-codedocstech_notesregex.md}

\subsection{\textasciitilde/code/docs/tech\_notes/001\_grep\_regex\_P\_examples.qmd}\label{codedocstech_notes001_grep_regex_p_examples.qmd}

\subsection{\textasciitilde/code/docs/tech\_notes/002\_grep\_examples.md}\label{codedocstech_notes002_grep_examples.md}

\subsection{SOME regex: in
\textasciitilde/code/zsh\_project/ZSH\_SH\_FILES/}\label{some-regex-in-codezsh_projectzsh_sh_files}

\subsection{REGEX}\label{regex}

TODO: - sed, when to use? - emphasize goal: use grep -P, regex to
understand how REGEX works. Tired of every 6 months learning all over
again. - greedy/not greedy and backtrack . Think like a regex engine! -
How to aerate regex ! - regex can be used to: - find - validate -
replace/insert - split - \ldots{} - When whiz, can do summersaults with
CLI, zsh tools (sed, grep , cut \ldots) and regex. Not NOW.

\begin{itemize}
\tightlist
\item
  Separate learning REGEX (grep -P, regex) and using REGEX in R, which I
  think is a tad easier.
\end{itemize}

\subsubsection{DEFINITIONS - as always,
crucial}\label{definitions---as-always-crucial}

\begin{itemize}
\item
  regex is a string; do not forget this.
\item
  META CHARACTERS - ascii (?) characters which by-default have
  non-literal meaning to engine that digests them. \textbf{Engine}
  specific. Must ESCAPE these characters to use as literals. Other
  contexts, such as unix shell, have similar idea: \texttt{\textless{}},
  \texttt{\textgreater{}\textquotesingle{}\ for\ example,\ refer\ to\ **redirect**\ .\ In\ C,\ sprintf,}\%`
  indicates formatting and literal use.
\item
  \textbf{To Escape} indicate to underlying engine that this meta
  character should be handled as though literal.
\item
  POSIX:
\item
  backslash\\
\item
  {[} {]}
\item
  \{ \}
\item
  ( )
\item
  caret \^{}
\item
  \$
\item
  dot .
\item
  pipe \textbar{}
\item
  ?
\item
  asterisk *
\item
  \texttt{+\ -}
\item
  ``+ - '\,'
\item
  \verbX + - X
\item
  \textbackslash begin\{verbatim\}
\item
  \begin{itemize}
  \tightlist
  \item
    \textbackslash end\{verbatim\}
  \end{itemize}
\end{itemize}

\subsubsection{Render REGEX Verbatim - 4 ways
(latex?)}\label{render-regex-verbatim---4-ways-latex}

\texttt{+\ -}

``+ - '\,'

\verb; + - ;

\verb;+ -;

\begin{verbatim} 
+ - 
\end{verbatim}

\begin{itemize}
\tightlist
\item
  \textbf{Character Class} Things like {[}0-9{]}.\\
  Rmk: {[}0-9{]}+ means repeat one or more of the prior
  \textbf{Character class} So both 321 and 333 match this regex.
\end{itemize}

\subsubsection{Specific to vim/neovim}\label{specific-to-vimneovim}

\begin{itemize}
\item
  magic = \v no need to escape (wait till know what doing first) - well,
  um.
\item
  magic = \texttt{\textbackslash{}v} no need to escape (wait till know
  what doing first)
\end{itemize}

\subsubsection{Specific to R}\label{specific-to-r}

\begin{itemize}
\tightlist
\item
  Before regex library (engine) sees code, the \textbf{compiler} (byte
  code?) gets it first. Must use double backslash for just one backslash
  to be seen by regex engine. Shell interpreters have no such compiler
  and single backslash suffices.
\end{itemize}

\#\#\#= Regex grouping: capture \& non-capture

\textbf{Perl} PCRE for lookaheads, capture (in R, perl=T) from !so

Groups that capture you can use later on in the regex to match OR you
can use them in the replacement part of the regex. Making a
non-capturing group simply exempts that group from being used for either
of these reasons.

Non-capturing groups are great if you are trying to capture many
different things and there are some groups you don't want to capture.

\subsubsection{LINUX / DEBIAN}\label{linux-debian}

RESCUE MODE, How to get to cli prompt? - Use Debian installer, rescue
mode - eventually will get prompt (a lot of screens to go through)

DEBIAN iso: should not need dd command; just cd to USB

\begin{itemize}
\item
  Debian: shutdown when sytem hangs? CTRL-ALT-DEL or ALT-printScreen
\item
  \emph{gnome-control-center} is cli to manage linux, X11, ..(even with
  cinnamon/lightdm)
\item
  xfce4-session-* \# (may need to poke around)
\item
  autologin, if using lightdm, as root edit /etc/lightdm/lightdm.confg

  \begin{itemize}
  \tightlist
  \item
    in {[}seat:*{]}, uncomment user and timeout line; set user to
    ``jim'' (no spaces around =)
  \end{itemize}
\item
  mount usb: sudo mount /dev/sdb /media/usb\_drive (RHS must EXIST)
\end{itemize}

versions: - Debian stable - Debian testing/sid - next version

\subsubsection{mount}\label{mount}

mount --bind mount \textbar{} grep ``/dev/'' \# list devices /dev/*

To temporarily load and run and sys: sudo chroot /mnt/X/ \# where /mnt/X
must already exist

\#\#\#= plymouth (linux) Bundled in initial ramdisk, before any fs
mounted. Provides a graphical boot?? splash screen check
/var/log/boot.log

\textless2024-10-04 Fri\textgreater{} Cinnamon and i3 do not mix well
(at this point) - Both are ``desktop environment'' (others: xfce4,
gnome?) - When lightdm (login manager) runs, user can choose 'desktop
envir'' AND authenticates. - SEE for ``rules''
https://forums.linuxmint.com/viewtopic.php?p=2533879\#p2533879 - LMDE is
meant to run Cinnamon (and not i3); I got i3 to run using testing
Debian. But this breaks LMDE (Cinnamon).

21-July-2023: Rumors, Linux Mint (now based upon Ubuntu/Canonical ) may
be moving to \textbf{LMDE} (Linux Mint Debian Edition) - 2024 - tried
Mint; I think had some issues with i3; tried Debian - 2025 - Debian 13

\begin{itemize}
\item
  Booting \ldots{} firmware \textbar{} bootloader (finds all kernels,
  os) \textbar{} grub2 (user select) ;\\
\item
  READ \textgreater info grub
\item
  /kernel is MINIMUM to start; this is why drivers often need separate
  install, not in kernel.
\item
  SWAP - latest Linuxmint built-in, no need
\item
  LinuxMint - installer sets mountpoints
\item
  PARTIONs - \home is separate; ~for all else (\textasciitilde30-40GB
  enough)
\item
  Boot drive - needs flag \texttt{boot} and \texttt{esp} (?)
\end{itemize}

\#\#\#= GRUB2 - 1st: do not panic - pre-linux boot - Usally all it wants
is to know which partition is root - Read: any simple article on this.
(ex: https://www.dedoimedo.com/computers/grub-2.html) - May see (hd0) or
(hd0,gpt2), (hd0,gpt3) refer to hard drive and partition - In my case
(hd0, 2), and command \textasciitilde{} set root=(hd0,gpt2) - May
involve updates to /etc/default/grub (use their tool)

GRUB is bootloade, boots kernel into RAM (vmlinuz, .initrd.gz) MBR long
gone, IGNORE (hd0,2) means drive hd0, partition 2

*** SETUP/CONFIG new HP Elitebook machine \textless2024-10-10
Thu\textgreater{} HP Elitebook but with pure Debian/Cinnamon \& i3

\emph{CAPSLOCK ESCAPE}; many easy ways =\textgreater{} confusion! -
Debian/Cinnamon \& i3: \emph{setxkbmap -option ``caps:escape''
(immediate)} reboot? may need to source \textasciitilde/.xinitrc - other
methods, like /etc/default/keyboard do not seem to work

(did not work for non-Cinnamon)CAPS\_LOCK: Use Cinnamon, keyboard
options GUI to set CAPS\_LOCK to ESC (easy) June 9, 2024 (Configure HP
Elitebook, 845, G8 - 2nd HP laptop) -
\emph{\textasciitilde/dotfiles/create\_soft\_links.sh} VERY helpful -
Emacs: must link emacs files in \textasciitilde/dotfiles files in
\textasciitilde/emacs.d* = - must re-install: wezterm (see webpage), i3,
zsh, git, emacs, ZSH, gh - public keys .. read debian's guide - zsh
change shell: chsh -s \$(which zsh) jim - emacs: FIRST install/config
\emph{`use-package'} always a pane: - Cinnamen:
hardward/keyboard/layout/options choice to set CAPSLOCK to ESC (works,
but with i3 too?) - websites: Google 1st, then Firefox (rest should
follow) - \emph{keep} sh files, links, config files up-to-date AND in
dotfiles/backup!

Drive is NOT 512G NAME MAJ:MIN RM SIZE RO TYPE MOUNTPOINTS sda 8:0 0
238.5G 0 disk ├─sda1 8:1 0 512M 0 part /boot/efi ├─sda2 8:2 0 27.9G 0
part / ├─sda3 8:3 0 977M 0 part {[}SWAP{]} └─sda4 8:4 0 209.1G 0 part
/home

\textbf{du -f (run on 11NOV2024)} Filesystem Size Used Avail Use\%
Mounted on udev 7.5G 0 7.5G 0\% /dev tmpfs 1.5G 1.7M 1.5G 1\% /run
/dev/sda2 28G 9.4G 17G 37\% / tmpfs 7.5G 18M 7.5G 1\% /dev/shm tmpfs
5.0M 12K 5.0M 1\% /run/lock /dev/sda4 205G 3.9G 191G 2\% /home /dev/sda1
511M 5.9M 506M 2\% /boot/efi tmpfs 1.5G 52K 1.5G 1\% /run/user/1000

*** format fat32, for copiers To make fat32 usb device.(for copier) 1)
(optional; too slow) can put all zeros (optional, slow -
\textasciitilde{} 5MB/s or 200 seconds for 1 GB on USB 2.0) 2) sudo
parted /dev/sda mklabel msdos (makes empty partition table, of form MBR)
3) sudo parted /dev/sda mkpart primary fat32 0\% 100\% (makes partition)
4) sudo mkfs.fat -F32 /dev/sda (format this partiion, if says to use -I
do so!) 5) sudo parted /dev/sdb print (confirms fat32)

To burn iso on usb (I never got gui's to work) 1) sudo dd bs=4M
if=/path/to/file.iso of=/dev/sda status=progress oflag=sync

\begin{itemize}
\tightlist
\item
  mount usb\_device (thumb drive)
\item
  format, partition etc. check fs: \textbf{df -Th /dev/sda}
\end{itemize}

-format -NO! use fat32 (above) to work with other devices sudo
mkfs.exfat /dev/sda

-mount mount /dev/sda /media/a\_mt\_point (a\_mt\_point must already
exist)

*** remap capslock to escape

\{ \# PURPOSE: \textbf{maps ChromeBox ``capslock'' key to Escape.} \# -
use \textgreater{} xev to find that capslock is key 133. \# - xmodmap is
older, but simpler to change key action to change key action. \# - newer
is \textbf{setxkbmap} but I find more effort to figure out simple
things. \# - SEE tech\_notes \# - lots of ways to do this remap. This
works, stay with it: \# DEPRECTED:

xmodmap -e ``keycode 133 = Escape'' Lenovo: capslock keycode = 66, and
escape is 9. However, capslock still insisting on going in caps lock
(UPPER CASE) setkbmap seems to suggest using caps:swapescape and not
caps:escape, but xmodmap won't accept.

\}

** cron job, crontab

\{ grep jim /var/log/syslog \# see cron jobs that ran

Sat May 21 18:48:16 PDT 2022 - jr changed /etc/rsyslog/50-default.conf -
uncomment \#cron -- cron s/d now log to cron.log - after change, run
sudo service rsyslog restart

\begin{itemize}
\tightlist
\item
  see cron Icard (`linux') \}
\end{itemize}

** Linux Kernel

\{ - one LTS Ubuntu can have many (upstream) kernels - Mix \& Match
kernels? X? - Kernel Upgrade - See INDEX C.

\}

\begin{itemize}
\tightlist
\item
  systemctl - control dameon and services
\item
  loginctl - control login daemon (such as when i3/lightdm crashes)
\end{itemize}

\emph{jim\_Permissions} u g o (user group other)

** ZSH notes

\emph{ZSH} ZSH REF CARD: (start)
https://www.bash2zsh.com/zsh\_refcard/refcard.pdf SEE MANUAL:
https://zsh.sourceforge.io/Doc/Release/ SEE ZSH GUIDE (2003, Stephenson)
https://zsh.sourceforge.io/Guide/zshguide.html ZSH FAQ (2010)
https://zsh.sourceforge.io/FAQ/ ZSH REF: Lots of simple GLOB examples
here:
http://reasoniamhere.com/2014/01/11/outrageously-useful-tips-to-master-your-z-shell/

\textasciitilde/dotfiles/.zshrc \textasciitilde/dotfiles/.zshenv zsh -x
*.sh \# prints line then executes

\subsubsection{script header}\label{script-header}

\#!/bin/zsh -xv \# verbose

\subsubsection{diff -y file1 file2 (2
columns)}\label{diff--y-file1-file2-2-columns}

\begin{itemize}
\tightlist
\item
  SEE github, zsh\_project
\item
  arrays, SEE 008\_array, ordered, unordered, append, typeset, unset,
  array + glob
\item
  structures: for, case \ldots{} see 032\_
\item
  find: 028\_ and others
\item
  prompt, ansi, echo \$fg\_bold{[}green'{]}hi SEE:
\item
  also: grep, \$\{var\}, tee, (( string\_to\_numbers)), do..done
\end{itemize}

\subsubsection{widgets}\label{widgets}

\begin{itemize}
\tightlist
\item
  software chunk, ``beginning of line'', ``kill-line''
\item
  to invoke: user types , \textless CTRL-{[}-B\textgreater{} (using vim
  notation)
\item
  but notation: `\^{}a', `\eb' `\^{}{[}b'
\item
  `\e' =
\item
  `\^{}{[}' = also
\end{itemize}

\subsubsection{bind to widget}\label{bind-to-widget}

\begin{itemize}
\tightlist
\item
  bindkey `\^{}k' kill-line (binkey.shortcut.existing widget)
\end{itemize}

\subsubsection{glob vs extended-glob (SEE Guide
5.9.7):w}\label{glob-vs-extended-glob-see-guide-5.9.7w}

\begin{itemize}
\tightlist
\item
  ll demo/**/*.txt
\item
  ls -l demo/**/*\textless1-10\textgreater.txt
\item
  ls -l demo/**/{[}a{]}*.txt (begins with \texttt{a})
\item
  ls -l demo/**/{[}\^{}cC{]}*.txt (do not begin with \texttt{c} or
  \texttt{C})
\item
  print -l **/*(/) \# just dir (recurse)
\item
  print -l **/*(.) \# just files (recurse)
\item
  ls -ldh demo/**/*(DOL{[}1,6{]}) \# 6 largest files
\item
  Lots of syntax to select JUST what you want!
\end{itemize}

\subsubsection{parameter expansion}\label{parameter-expansion}

SEE 000\_PE\_examples.sh - \$\{var/old/new\} 1 replace -
\$\{var//old/new\} all - \$\{f\#\emph{.\} is x.y where f = file.x.y
(remove left) - \$\{f\#}.\} is x.y where f = file.x.y (remove left) -
\$\{f\#\#*.\} is y (greedy, remove all left)

\begin{itemize}
\tightlist
\item
  SEE FAQ73, https://mywiki.wooledge.org/ , GNU Bash Ref Manual 3.5.3
\end{itemize}

\subsubsection{zsh/unix permissions}\label{zshunix-permissions}

*** unix, zsh permissions u g o owner--in group--others

r=4 w=2 x=1 chmod 700 rwx --- --- ``public'' 644 rw- r-- r-- ``private''
600 rw- --- ---

*** zsh commands in emacs (SEE: zsh\_project)

\subsubsection{completion, Use zstyle}\label{completion-use-zstyle}

man zshcompsys zstyle + 3rd party scripts - do MANY things (completion,
modify settings, config VCS\_INFO \ldots) SEE:
https://thevaluable.dev/zsh-completion-guide-examples/

\emph{VCS\_INFO} is a function, used to populate variables (prompt, for
ex) \emph{retrieved from vcs, ie git} SEE:
https://arjanvandergaag.nl/blog/customize-zsh-prompt-with-vcs-info.html
SEE:
https://zsh.sourceforge.io/Doc/Release/User-Contributions.html\#Version-Control-Information

USAGE: \emph{zstyle}

and is NOT so OBVIOUS ! if match, apply the style USAGE: zstyle

Completion: ex: cd completion zstyle `:completion:\emph{:}:cp:' zstyle
':completion:*'

General: :completion::::: - command (cd, rm, etc) - (could be files,
users, options ??)

\emph{precmd} is hook function, runs before ZSH prompt (SEE manual
9.3.1) SEE:
https://zsh.sourceforge.io/Doc/Release/Functions.html\#Special-Functions

\subsubsection{printf}\label{printf}

\begin{itemize}
\item
  string(\%s), digit(\%d); example: digit,pad with by 5 printf
  ``\%05d\n'' \$\{x\}
\item
  R: see sprintf https://www.r-bloggers.com/2010/05/number-formatting/
\end{itemize}

*** zle = zsh line editor/keymaps/widgets/

*** misc zsh

SEE Serge Gebhardt widgets correspond to commands, often with shortcut

Given a widget (ex: `\eb', ie esc b) bindkey `\eb' \#returns
backward-word; ie ESC-b, aka M-b goes back one word !\\
binddky `\^{}b' \# returns backward-charac

bindkey `\ef' \# returns foward-word

bindkey `\^{}a' \# returns beginning-of-line

zle -al \# list of widgets

keymaps=collection of shortcuts bindkey -l \# vi modes?

To see mapping: 1. cat ; type key; ends (try M-b) 2. ; type key

invoke widget: 1. bind to shortcut 2. zle \# to execut 3. sp widget

TERMCAP - obsolete

\subsubsection{sudo vs su \ldots.}\label{sudo-vs-su-.}

\{

\begin{itemize}
\item
  su jim change to User `jim'
\item
  sudo cmd - last \textasciitilde{} 15' (temporary use of root
  privileges) - asks for user's password - allows root `privileges' but
  the home directory, path etc remains the user's
\item
  next 3 are same: (change to root user, with root env) {[}use exit to
  exit{]}
\item
  \begin{itemize}
  \tightlist
  \item
    su - root \# change to root user, with root env
  \end{itemize}
\item
  \begin{itemize}
  \tightlist
  \item
    sudo su
  \end{itemize}
\item
  \begin{itemize}
  \item
    sudo su -

    \begin{itemize}
    \tightlist
    \item
      shell: either login or non-login
    \item
      non-login has 2 flavors: \textbf{interactive} (user at CLI) and
      \textbf{non-interactive} (a subshell for scripts)
    \end{itemize}
  \end{itemize}
\end{itemize}

!askubuntu 376199 !askubuntu 1225041

\}

\subsubsection{drive info}\label{drive-info}

\begin{verbatim}
# succinct, useful info
lsblk --output NAME,UUID,PARTUUID
\end{verbatim}

xev keyboard

\{ - Keyboard specific, find what \emph{keycode} a button is mapped to:
- USAGE: \textgreater{} xev - type just 1 button, look for its keycode,
keysym on this keyboard - example: q will be keycode=24, keysym=0x71
called `q'

\}

17JULY2023 - ebook-viewer (calibre) has conflict with caps:swapescape,
can not figure out REF:
\url{https://wiki.archlinux.org/title/Xorg/Keyboard_configuration} FIX:
now using \textbf{setxkbmap -option caps:escape (in .xinitrc) }
DEPRECATES anything before 17JULY 2023

xxd

\{

\begin{verbatim}
- To find how zsh maps a button (A, alt, F2) :  
- USAGE: > xxd <CR>
- press <ALT>+a
- terminal displays coding (^[a)
- SEE ROTHGAR
\end{verbatim}

\}

\subsubsection{GLOB}\label{glob}

grep\_vs\_ls

\emph{Grep} always finds words that match a pattern and returns file
names of matches.

ls (+ glob) \textbf{finds filenames that match a pattern}. Very
different from grep. (same in vim)

\emph{jim\_GLOB\_examples} Mostly of form ls or ll or print -l and **/*
example: print -l \textasciitilde/code/**/*.(R\textbar Rmd) \# any
level, return all .R and .Rmd files

See my zsh GLOG handwritten notes (till typed in here)

\subsubsection{ZLE}\label{zle}

ZLE = Zsh line editor \textbar{} NOT GNU readline\\
\emph{zle\_widgets} (all commands)

Output from zle -al (\textasciitilde403 cmds)

\subsubsection{BINDKEY}\label{bindkey}

\emph{bindkey} \# results, all shortcuts

(sample) ``\textsuperscript{A''-''}C'' self-insert ``\^{}D''
list-choices ``\textsuperscript{E''-''}F'' self-insert ``\^{}G''
list-expand ``\^{}H'' vi-backward-delete-char ``\^{}I''
expand-or-complete

\subsubsection{Google, API, curl}\label{google-api-curl}

(1APR2022) Google's example, with loop for uri\_redirect
https://accounts.google.com/o/oauth2/v2/auth?
scope=https\%3A\%2F\%2Fwww.googleapis.com\%2Fauth\%2Fyoutube.readonly\&
response\_type=code\&
state=security\_token\%3D138r5719ru3e1\%26url\%3Dhttps\%3A\%2F\%2Foauth2.example.com\%2Ftoken\&
redirect\_uri=http\%3A//127.0.0.1\%3A9004\& client\_id=client\_id

\begin{verbatim}
- Google's authorization server: https://accounts.google.com/o/oauth2/v2/auth
\end{verbatim}

\subsubsection{From Explorer}\label{from-explorer}

GET
https://youtube.googleapis.com/youtube/v3/playlists?part=snippet\%2CcontentDetails\&maxResults=5\&mine=true\&key={[}YOUR\_API\_KEY{]}
HTTP/1.1

Authorization: Bearer {[}YOUR\_ACCESS\_TOKEN{]} Accept: application/json

\subsubsection{}\label{section-2}

same, but as Curl

curl\\
`https://youtube.googleapis.com/youtube/v3/playlists?part=snippet\%2CcontentDetails\&maxResults=5\&mine=true\&key={[}YOUR\_API\_KEY{]}'\\
--header `Authorization: Bearer {[}YOUR\_ACCESS\_TOKEN{]}'\\
--header `Accept: application/json'\\
--compressed

\subsubsection{From Google Playground}\label{from-google-playground}

https://youtube.googleapis.com/youtube/v3/commentThreads?videoId=Mec9sw1cJk8\&part=snippet,replies
\#\#\#

\subsubsection{}\label{section-3}

CURL \textbar{} YOUTUBE API \textbar{} GOOGLE API \textbar{} OAUTH 2.0
\textbar{}

client = oauth\_client(id= client\_id, token\_url = token\_url, secret =
client\_secret, key = API\_KEY, auth = ``body'', \# header or body

\begin{verbatim}
    name = "youtube_ONE_video_ALL_comments")
\end{verbatim}

req \textless-
request(``https://www.googleapis.com/youtube/v3/commentThreads?videoId=Mec9sw1cJk8\&part=snippet,replies'')
\%\textgreater\% req\_oauth\_auth\_code(client = client, auth\_url =
auth\_url, token\_params=scope{[}{[}1{]}{]})

resp \textless- req \%\textgreater\% req\_perform()

Some Remarks: - Google is but one implementation of various API, oauth
technologies. The more you read the more confused you may become (at
least for me).\\
- The R package \textbf{gargle} is uses \textbf{httr} and therefore not
my preference.\\
- I am using httr2 to automate things; I'd like to understand things
using a little as possible: curl, browser and local server running as
localhost.\\
- Most of the R work is done at lower level, such as packages curl and
httpuv.

begin\{verbatim\} G O O G L E end\{verbatim\}

\subsubsection{HTTR2 - NOTES (needs clean
up!)}\label{httr2---notes-needs-clean-up}

PURPOSE: Demonstrate configuration for HTTR2 and OAUTH2 with Google's
Youtube API.

\begin{verbatim}
                        - uses off-the-shelf `httr2::req_oauth_auth_code()` + configuration
                      - uses authorization code flow.
                        - uses redirect_uri localhost, cut & paste (via obo) is deprecated.
                        - httr2:: hides almost all details of interaction.
                        - use curl and localhost such as httpuv:: to see lower level
\end{verbatim}

Source:
https://developers.google.com/youtube/v3/guides/auth/installed-apps

RELATED INFO: - Google Explorer (youtube) - Google OAUTH2 playground

\begin{verbatim}

#\###**==========================
From Google (Youtube) Explorer:
GET https://youtube.googleapis.com/youtube/v3/playlists?part=snippet%2CcontentDetails&maxResults=5&mine=true&key=[YOUR_API_KEY] HTTP/1.1

Authorization: Bearer [YOUR_ACCESS_TOKEN]
Accept: application/json

#\###**==========================
\end{verbatim}

For youtube (auth code): echo ``curl -Lsv
"https://accounts.google.com/o/oauth2/v2/auth?\\
client\_id=\$client\_id\&\\
redirect\_uri=https://127.0.0.1:8080\&\\
scope=https://www.googleapis.com/auth/youtube\&\\
response\_type=code"''

scope = list( ``https://www.googleapis.com/auth/youtube'',
``https://www.googleapis.com/auth/youtube.force-ssl'')

For youtube (obtain results): curl\\
`https://youtube.googleapis.com/youtube/v3/playlists?part=snippet\%2CcontentDetails\&maxResults=5\&mine=true\&key={[}YOUR\_API\_KEY{]}'\\
--header `Authorization: Bearer {[}YOUR\_ACCESS\_TOKEN{]}'\\
--header `Accept: application/json'\\
--compressed

\#\#\#= NEEDED SCOPES: https://www.googleapis.com/auth/youtube Manage
your YouTube account https://www.googleapis.com/auth/youtube.force-ssl
See, edit, and permanently delete your YouTube videos, ratings, comments
and captions

playlistId = ``PLlXfTHzgMRUIqYrutsFXCOmiqKUgOgGJ5'' \# Pavel Grinfeld,
Linear Alg 3

\begin{verbatim}
                E N D G O O G L E
\end{verbatim}

\#\#\#= Procedure: - Follow hadley outlines in Vignette for Github and
and getting user's information. (Requires oauth token) - Change for
google - let httr2 handle the details, use this function:
httr2::req\_oauth\_auth\_code() - If I have this right, this will (1)
get the access token and (2) complete REST request.

** Pandoc PANDOC: !pandoc --metadata=project:xxx --lua-filter
doc/panvimdoc/scripts/skip-blocks.lua --lua-filter
doc/panvimdoc/scripts/include-files.lua -t
doc/panvimdoc/scripts/panvimdoc.lua \% -o doc/jimHelp.txt

** MORE CURL CURL Examples:

cURL write (to standard) w response after callling example.com

\begin{verbatim}
curl -w "Response %{response_code}\n" example.com
\end{verbatim}

github curl https://api.github.com/zen

returns lot of kev=value pairs curl https://api.github.com/users/defunkt

-include headers curl -i https://api.github.com/users/defunkt

headers only curl --head

CURL\_CONFIG (a FILE) USAGE curl -K CURL\_CONFIG \ldots{}

\begin{verbatim}
url = example.com
-w "Type: Hello %{local_ip} \n"
\end{verbatim}

Misc Notes: ``State'' - cookies used to be used; now state carried in
headers

Misc Notes: ``State'' - cookies used to be used; now state carried in
headers\\
vim:nospell

** ANDROID

\subsubsection{Android, Mobile, Cell Phone -
notes}\label{android-mobile-cell-phone---notes}

\emph{Nov 2024} - need wireless phone? or working wifi-only phone? (mine
does not work)

\emph{Carrier} - 5G or \texttt{slower} 4G/LTE - 1-5 GB per month seems
OK for some video, streaming \ldots{} - podcasts \textasciitilde{}
1MB/minute of talk (use wifi) - want unlimited text/talk - MintMobile
(owned by T-Mobile) claims \$15 per month - US - SIMS: physical or eSim,
if device supports - not my current onePlus 3T,

\begin{itemize}
\tightlist
\item
  device must support 5G? GSM? Volte? ..
\item
  device should also support wifi (for mp3 downloads)
\end{itemize}

\emph{mp3 players} - DAC, dongle, digital devices: seem to be superb
music; not android or podcast devices. - Sony as usual is best; also
Fifo and many others

RETAIL DEMO UNIT (`retail mode'): - ie runs in endless LOOP, no Cell
ability, no MEI - useable ONLY for wifi - Can be BARGAIN, but \ldots{} -
Must unlock bootloader (to remove `endless loop software' and become
regular wifi device. Locked means bootloder hard-coded insist OS match a
code. - if CAN unlock bootloader , BARGAIN. Beware endless hours
otherwise.

\begin{itemize}
\tightlist
\item
  bootloader lock - at least TWO kinds
\item
  carrier lock - that carrier only and even if unlocked software may
  persist
\item
  factory lock - seems many are carrier unlocked; but less or factory
  unlocked.
\end{itemize}

*** Google Pixel \textbf{3a XL} (my phone, 6.3'' good size) - screen
broken Aug '25 - Android 12 = final google update. - DO expect
``unofficial'' Android 13 for this phone (sooner or lalter).

\begin{itemize}
\item
  Pixel 4a (5.8'', too small) cheaper glass than pixel 3a; practical
  gone (Aug '25)
\item
  OnePlus
\item
  (1/26) - older ``flagships'' models should be good.
\item
  Samsung - some say good; other don't agree.
\item
  waste so much time on this.
\item
  OnePlus 12, 13 all good - not cheap
\item
  OnePlus 8 (bought 9/25) (\$137) very good, but no 3mm jack;
\item
  buy Redmi Note 14; because mid-range; has 3mm jack; 6GB/128GB (\$140)
\end{itemize}

** Typst (and emacs) (22 MARCH 2026) - assorted problems getting all
emacs pieces installed. needs lsp, mist, and a few other pieces. Maybe
very active development now. Postpone till I have more info.

*** \textbf{OEM Unlock} - greyed out? (like mine) then not possible to
unlock bootloader itself. Means: no root. \textbf{no ROM install}
\textbf{no TWRP} - my pixel is VERISON (sprint?) phone; not a Google
phone; b/c IMEI begins with 35\ldots{} NO way to change bootloader.

\subsubsection{ADB}\label{adb}

\begin{itemize}
\item
  \emph{ADB} DEBUG: a ``mode'' that allows installing apps, read logs on
  Android, file transfer\ldots{} Works by running TCP sever on host (PC)
  and daemon on device (phone) Works by running TCP sever on host (PC)
  and daemon on device (phone)
\item
  SEE
  https://www.howtogeek.com/192732/android-usb-connections-explained-mtp-ptp-and-usb-mass-storage/

  \begin{itemize}
  \tightlist
  \item
    adb --help
  \item
    adb shell
  \item
    adb shell df -h

    \begin{itemize}
    \tightlist
    \item
      adb shell ls /system/bin (available cmds)
    \end{itemize}
  \item
    adb push {[} -- sync {]} \# only push if newer
  \item
    \textbf{adb push * /sdcard/Music \# worked 3/25, push to pixel 4}
  \end{itemize}
\item
  Photo tranfer, different.
\item
  \textbf{mp3 file transfer} ADB appears to be FUSSY: remove things like
  \texttt{?} from file names or foregin char. ADB sucks at error
  messages; chokes; just seem to stall. Just fix the file names and adb
  will work; speed is very good; but even 25 MB/s \textasciitilde{} 1.8
  GB/min. Be patient with 40 GB.
\item
  Bluetooth - wasted plenty of time: use wired ADB; some mention Ubuntu
  \& bluetooth never got along. Either way - TIME SINK; waste.
\item
  \emph{MTP} is protocol to move files; seems imperfect (CLAIM: now
  standardized, better) ; PTP for photos
\end{itemize}

mtp://{[}usb:001,085{]} where 085 refers to device. (Run lsusb) - AVOID
this stuff; \textbf{stay with ADB and fix those file names}

\begin{itemize}
\tightlist
\item
  \textbf{adb backup} dissappointing, time-sink; THINK backup all all my
  apps, data, but can NOT find clear documentation. STOP.
\item
  Do not like Google bloatware. Expected something like ONEPLUS (which I
  install ROM). Google's rules, annoyances - must remove. do not want G-
  ecosystem to point to each of its sister apps.
\end{itemize}

*** Android os - Android is U/I to actual OS, which is \textbf{Dalik},
uses java VM \textbf{Recovery Mode} is separate partition(?) contains
just enough code to boot in this mode. Replacing this code is
\textbf{custom} recovery vs \textbf{stock} recovery. - \textbf{FASTBOOT}
purpose to \texttt{flash} ROM on device; level beyond ADB.

\subsubsection{section\{Laptop Buying
Notes\}}\label{sectionlaptop-buying-notes}

\textbf{eMMc} is on bmotherboard(embedded), slow but works: cheap,
reliable; fine to boot. Check /dev/mmcblk** SSD is much better, but more
expensive.

*** Lenovo T480 (stolen May 2024)

\begin{itemize}
\tightlist
\item
  running Linux Mint (no more Chromebooks)
\item
  Power, 65W, need brick or wall charger.
\item
  cable must support 3 Amps
\item
  buzz words GaN, PD (Power Delivery), no need latest PPS
\item
  name brands: Anker, Belkin, ``Amazon Basics'', \textbf{Beware
  off-brand} buy based on what is compatible with T480 (go crazy trying
  to match standards, USB-C 3.1, 3.2, 4.0, generatations, standard or
  not?)
\end{itemize}

** RUST (systems level language)

\begin{itemize}
\item
  programmer has control over memory, variables. Leads to SAFETY and
  PERFORMANCE. At cost of understanding more about memory etc.
\item
  \emph{macro} is code that runs at COMPILE time; inserts compiled code
  sections. (Saves programmer from needing to write common code over and
  over.)
\end{itemize}

\emph{CAR NOTES} https://check-vin.org/

my VIN: 1HGEM22911L031079

\subsubsection{Battery Notes (3/1/2026)}\label{battery-notes-312026}

\begin{itemize}
\tightlist
\item
  AGM is lead, with sulfuric acid. The acid is absorbed(?) by a mat; so
  does not spill.
\item
  Resting Voltage (no load) \textasciitilde{} 12.8-12.9 V
\item
  When just recharged - voltage can be \textgreater{} 13 V. This will
  settle back soon.
\item
  When car starts, the 12.8V will drop (should stay above 10V) and
  immediately bounce back as alternator begins.
\item
  The alternator (regulates battery re-charge) will bring voltage up to
  \textasciitilde14V, at least briefly.
\end{itemize}

Because computer runs the show, any change in one place triggers change
somewhere else computer constantly adjusts.

DC Voltage 100\% charge - 12.8 75\% charge - 12.6 50\% charge - 12.3

Time to recharge a 12V battery depends on SoC (state of charge?) and amp
output (10, 20 even 30 AMPS) of charger. Noco is good name. Battery
charger is NOT same as battery jumper, the latter is more expensive. Ask
because I don't know.

BLACK, Prius 5 - 4 door , all owners in family (?) Chicago car:
https://corvallis.craigslist.org/cto/d/corvallis-2011-toyota-prius-4dr/7791577925.html
carfax:
https://www.autovhr.com/CFReports-ran/AutoVhr-CFreport\_1728601520.html
-carfax clean no records, post-2018 Ingrid 2009/Prius Hatch
4D/112k/\$9/(carfax: says \$8.4) (JTDKB20U3 93503700) car: WHTIE,
https://portland.craigslist.org/wsc/cto/d/portland-upgraded-base-model-prius-2009/7786557907.html
carfax:
https://www.autovhr.com/CFReports-ran/AutoVhr-CFreport\_1728671601.html

waiting: Travis has receipt for last; and battery: 12V or hybrid

\#------------------------ Jack JACK,
2012/preisC/111K/\$8.9/JTDKDTB30C1023315 (\emph{dealer offered him \$7k
trade-in}) Waiting: his battery, recent records car:
https://portland.craigslist.org/mlt/cto/d/portland-2012-toyota-prius/7789997037.html
carfax: see pdf (shows battery replaced; Jack says it wasn't) service
records: see pdf

\subsubsection{Reddit, API}\label{reddit-api}

\begin{verbatim}
Managing and exporting your saved items on Reddit can be tricky because **Reddit’s native interface does not include a search bar for saved items**, and it enforces a strict **1,000-item limit** (saving a 1,001st item will push the oldest one out of your list).

Because your saved list is private, a direct URL like `reddit.com/user/jimrothstein1/saved/` can only be viewed when you are securely logged into your specific account.

A breakdown of the different methods available to search, download, or handle this list includes Python options, AI tools, and browser extensions.


### Method 1: Using Python and the Reddit API (Highly Recommended)

Using Python alongside **PRAW** (Python Reddit API Wrapper) is the most reliable way to back up your list. It allows you to download everything into a readable file (like a CSV or Markdown file) which you can then easily search using standard computer tools or open in Excel.

**How it works:**

1. Go to Reddit’s App Preferences while logged in and create an app to get your `client_id` and `client_secret`.
2. Install PRAW using pip: `pip install praw`
3. Run a script to iterate through your saved history.

**Example Python Script:**

```python
import praw
import csv

# Initialize Reddit API credentials
reddit = praw.Reddit(
    client_id="YOUR_CLIENT_ID",
    client_secret="YOUR_CLIENT_SECRET",
    user_agent="SavedPostExporter v1.0 by /u/jimrothstein1",
    username="jimrothstein1",
    password="YOUR_PASSWORD"
)

# Open a CSV file to write the data
with open("reddit_saved.csv", "w", newline="", encoding="utf-8") as file:
    writer = csv.writer(file)
    writer.writerow(["Type", "Title/Comment", "Subreddit", "URL"])

    # Fetch saved items (Maxes out at Reddit's 1000 limit)
    for item in reddit.user.me().saved(limit=None):
        if isinstance(item, praw.models.Submission):
            writer.writerow(["Post", item.title, item.subreddit.display_name, item.url])
        elif isinstance(item, praw.models.Comment):
            writer.writerow(["Comment", item.body, item.subreddit.display_name, item.permalink])

print("Backup complete! Check reddit_saved.csv")
\end{verbatim}

\#---

\subsubsection{Method 2: Can you tell an AI Agent to do
it?}\label{method-2-can-you-tell-an-ai-agent-to-do-it}

\textbf{Yes, but with security and privacy caveats.}

An AI agent \textbf{cannot} download the list simply by reading the URL
you provided. Because your saved list requires a secure login, a
standard AI agent will just see a ``403 Forbidden'' or ``Page Not
Found'' error.

To use an AI agent (like a local LLM or an AI automation tool), you
would have to choose one of two paths:

\begin{enumerate}
\def\labelenumi{\arabic{enumi}.}
\tightlist
\item
  \textbf{The Automated Browser Route (Safer but Technical):} You can
  use an AI coding assistant (like Claude or ChatGPT) to write a script
  utilizing \textbf{Playwright} or \textbf{Selenium}. You log into
  Reddit manually on your computer, and the AI-generated script handles
  scrolling through the page, scraping the text, and compiling it into a
  document.
\item
  \textbf{The Credential Route (High Risk):} You could provide an AI
  agent with your Reddit API keys or username/password to execute a
  Python backup script for you. \textbf{This is highly discouraged}
  unless the AI agent is running completely locally on your own machine,
  as sharing passwords or session cookies with third-party cloud AIs
  poses a severe security risk.
\end{enumerate}

\subsubsection{Method 3: Dedicated Third-Party Tools \&
Extensions}\label{method-3-dedicated-third-party-tools-extensions}

If you prefer not to write code, several community-built tools bypass
Reddit's limitations:

\begin{itemize}
\tightlist
\item
  \textbf{Reddit Data Request (Official Takeout):} You can request a
  complete legal archive of your data directly from Reddit via
  \textbf{\href{https://www.google.com/search?q=https://reddit.com/settings/data-request}{reddit.com/settings/data-request}}.
  Reddit will email you a ZIP file containing CSV files of your history.
  Note that this often only contains the \emph{links} to your saved
  items rather than full text descriptions, and it can take several days
  to process.
\item
  \textbf{Browser Extensions (Notion / CSV Exporters):} There are
  several popular extensions on the Chrome Web Store and Firefox Add-ons
  library designed specifically for this purpose. They use secure OAuth
  authentication (meaning you never reveal your password) to let you
  click a single button to export your saved items directly into a
  spreadsheet or a Notion database, complete with search and tagging
  capabilities.
\item
  \textbf{Web Utilities (e.g., Reddit-Saved.com):} Free web apps exist
  that securely authenticate through Reddit OAuth, index your list, and
  provide an instant search bar. Because they store a record of your
  saves once synced, they can often help you preserve items beyond
  Reddit's 1,000-post memory wall.
\end{itemize}

```




\end{document}
